\documentclass[11pt,letterpaper]{article}

\usepackage[margin=1in]{geometry}
\usepackage{booktabs}
\usepackage{longtable}
\usepackage{array}
\usepackage{hyperref}
\usepackage{enumitem}
\usepackage{xcolor}
\usepackage{titlesec}
\usepackage{fancyhdr}
\usepackage{tabularx}

% Brand colors
\definecolor{garnet}{HTML}{73000A}
\definecolor{atlantic}{HTML}{466A9F}
\definecolor{congaree}{HTML}{1F414D}
\definecolor{rose}{HTML}{CC2E40}
\definecolor{black90}{HTML}{363636}

\hypersetup{
    colorlinks=true,
    linkcolor=garnet,
    urlcolor=atlantic,
    citecolor=congaree
}

\titleformat{\section}{\Large\bfseries\color{garnet}}{}{0em}{}[\titlerule]
\titleformat{\subsection}{\large\bfseries\color{congaree}}{}{0em}{}
\titleformat{\subsubsection}{\normalsize\bfseries\color{black90}}{}{0em}{}

\pagestyle{fancy}
\fancyhf{}
\fancyhead[L]{\small\color{black90}CSCE 775 -- Deep Reinforcement Learning}
\fancyhead[R]{\small\color{black90}Study Plan: Lectures 4--7}
\fancyfoot[C]{\thepage}
\renewcommand{\headrulewidth}{0.4pt}

\setlength{\parindent}{0pt}
\setlength{\parskip}{6pt}

\begin{document}

\begin{center}
{\LARGE\bfseries\color{garnet} CSCE 775: Deep Reinforcement Learning\\[4pt] Comprehensive Study Plan}\\[8pt]
{\large Lectures 4--7: Model-Free RL, Deep Learning, PyTorch, DAVI \& DQNs}\\[4pt]
{\normalsize J.C. Vaught \quad$\cdot$\quad Spring 2026 \quad$\cdot$\quad University of South Carolina}
\end{center}

\vspace{6pt}
\hrule height 1pt
\vspace{12pt}

\tableofcontents
\newpage

%=============================================================================
\section{Area 1: Model-Free Prediction (MC \& TD Methods)}
%=============================================================================

\subsection{Textbook Exercises (Sutton \& Barto, 2nd Ed.)}

\subsubsection{Chapter 5: Monte Carlo Methods}

Priority exercises: 5.1--5.6, and especially \textbf{5.12} (programming).

\begin{longtable}{@{}cp{4.8in}@{}}
\toprule
\textbf{Ex.} & \textbf{Description} \\
\midrule
\endhead
5.1 & Analyze why value function plots for blackjack show jumps at certain states \\
5.2 & Compare bias properties of every-visit and first-visit MC estimators \\
5.3 & Draw and describe the backup diagram for Monte Carlo estimation of $Q(s,a)$ \\
5.4 & Modify MC ES pseudocode to use incremental mean updates instead of storing all returns \\
5.5 & Calculate both first-visit and every-visit estimator types on a concrete example \\
5.6 & Derive importance sampling equations for action value estimates $Q(s,a)$ \\
5.7 & Explain the behavior of MSE in weighted importance sampling at the start of learning \\
5.8 & Determine variance properties of the every-visit MC estimator \\
5.9 & Modify the off-policy MC algorithm using incremental update formulas \\
5.10 & Derive the recursive weighted-average update rule \\
5.11 & Explain why the importance sampling ratio starts from $t+1$ in off-policy control \\
\textbf{5.12} & \textbf{Racetrack problem}: Implement MC control for a racetrack gridworld with noisy acceleration \\
\bottomrule
\end{longtable}

\subsubsection{Chapter 6: Temporal-Difference Learning}

Priority exercises for prediction: 6.1, 6.2, 6.3, 6.4, 6.5, 6.6.

\begin{longtable}{@{}cp{4.8in}@{}}
\toprule
\textbf{Ex.} & \textbf{Description} \\
\midrule
\endhead
6.1 & Derive the telescoping sum identity for TD errors (fundamental) \\
6.2 & The ``buffet'' thought experiment---explain when TD outperforms MC \\
6.3 & Compute value updates after the first episode of the 5-state random walk \\
6.4 & Analyze how different $\alpha$ values affect TD vs.\ MC performance \\
6.5 & Investigate how increasing learning rate affects MSE over multiple episodes \\
6.6 & Describe methods for computing the exact values in the random walk MRP \\
\bottomrule
\end{longtable}

\subsubsection{Chapter 7: $n$-step Bootstrapping}

Both exercises are essential. Exercise 7.2 gives hands-on experience with the bias-variance spectrum.

\begin{longtable}{@{}cp{4.8in}@{}}
\toprule
\textbf{Ex.} & \textbf{Description} \\
\midrule
\endhead
7.1 & Show the $n$-step return error can be written as a sum of TD errors \\
\textbf{7.2} & \textbf{Programming}: Compare $n$-step TD methods for different $n$ on the 19-state random walk, reproducing Figure 7.2 \\
\bottomrule
\end{longtable}

\subsection{University Assignments}

\begin{longtable}{@{}p{2in}p{3.8in}@{}}
\toprule
\textbf{Assignment} & \textbf{What You Do} \\
\midrule
\endhead
David Silver's Easy21 (UCL) & The gold standard. 5 parts: implement Easy21 env, MC control, SARSA($\lambda$), linear FA, discussion. \href{https://www.davidsilver.uk/wp-content/uploads/2020/03/Easy21-Johannes.pdf}{Assignment PDF} \\
\addlinespace
Georgia Tech CS 7642 Project 1 & Reproduce Sutton (1988) original TD experiments on bounded random walk. \href{https://omscs.gatech.edu/cs-7642-reinforcement-learning}{Course page} \\
\addlinespace
Stanford CS234 Module 3 & TD policy evaluation, batch MC vs batch TD. \href{https://web.stanford.edu/class/cs234/}{Course page} \\
\bottomrule
\end{longtable}

\subsection{Mini-Projects}

\subsubsection{Project A: MC vs TD(0) Random Walk Comparison}

Implement a configurable $N$-state random walk MRP (support $N=5, 19, 99$). Implement first-visit MC prediction and TD(0) prediction. Reproduce Figure 6.2 from Sutton and Barto: plot RMS error vs.\ episodes for TD(0) with $\alpha \in \{0.05, 0.1, 0.15\}$ and constant-$\alpha$ MC with $\alpha \in \{0.01, 0.02, 0.04\}$. Reproduce Example 6.3 (batch updating). Extend to 19 states and reproduce Figure 7.2.

\textit{Deliverables}: Four publication-quality figures, a written analysis (1--2 pages) of the bias-variance tradeoff observed.

\subsubsection{Project B: Blackjack Policy Evaluation Showdown}

Use Gym's Blackjack-v0 environment. Implement first-visit MC, every-visit MC, TD(0), and $n$-step TD prediction for the fixed policy: stick on 20 or 21, hit otherwise. Measure convergence speed. Produce 3D surface plots of the value function at 10K, 100K, and 500K episodes.

\subsection{Key Coding Resources}

\begin{itemize}[nosep]
    \item \href{https://github.com/dennybritz/reinforcement-learning}{Denny Britz RL Repository} -- MC/TD notebooks with starter code and solutions
    \item \href{https://github.com/ShangtongZhang/reinforcement-learning-an-introduction}{Shangtong Zhang Figure Reproductions} -- Python reproductions of nearly every textbook figure
    \item \href{https://github.com/udacity/deep-reinforcement-learning/blob/master/monte-carlo/Monte_Carlo.ipynb}{Udacity Deep RL MC Notebook}
\end{itemize}

\subsection{Recommended Sequence}

\begin{enumerate}[nosep]
    \item Read S\&B Ch 5.1--5.5, do exercises 5.1--5.6
    \item Code Denny Britz MC Prediction notebook (Blackjack)
    \item Read Ch 6.1--6.3, do exercises 6.1--6.6
    \item Code Mini-Project A (random walk, Figures 6.2 and 7.2)
    \item Read Ch 7.1--7.2, do exercises 7.1--7.2
    \item Complete Easy21 Parts 1--3
\end{enumerate}

\newpage
%=============================================================================
\section{Area 2: Model-Free Control (SARSA, Q-Learning, On/Off-Policy)}
%=============================================================================

\subsection{Textbook Exercises (Chapter 6)}

Priority: 6.6--6.7 (Windy Gridworld + extensions), 6.9 (Expected SARSA on cliff walking), 6.11--6.12 (why Q-learning is off-policy; greedy SARSA vs Q-learning), 6.14 (maximization bias / Double Q-learning).

\begin{longtable}{@{}cp{4.8in}@{}}
\toprule
\textbf{Ex.} & \textbf{Description} \\
\midrule
\endhead
6.6 & Implement SARSA on Windy Gridworld (Example 6.5); plot episodes vs time steps \\
6.7 & Extend 6.6 with King's moves (8 actions) and stochastic wind \\
6.8 & Derive the relationship between returns and action-value estimates through TD decomposition \\
6.9 & Programming: Compare SARSA, Q-learning, and Expected SARSA on cliff walking \\
6.10 & Why Q-learning's policy (argmax) differs from its behavior policy ($\varepsilon$-greedy) \\
6.11 & Explain precisely why Q-learning is classified as off-policy control \\
6.12 & ``If action selection is greedy, is Q-learning exactly the same as SARSA?'' \\
6.13 & Explore double learning applied to Expected SARSA \\
6.14 & Programming: Maximization bias problem, Q-learning vs Double Q-learning \\
\bottomrule
\end{longtable}

\subsection{University Assignments}

\begin{longtable}{@{}p{2in}p{3.8in}@{}}
\toprule
\textbf{Assignment} & \textbf{What You Do} \\
\midrule
\endhead
Easy21 Parts 3--5 & SARSA($\lambda$) with eligibility traces, linear FA comparison, discussion \\
\addlinespace
Stanford CS234 Asgn.\ 2 & Linear Q-learning then full DQN on Atari. \href{https://web.stanford.edu/class/cs234/}{Course page} \\
\addlinespace
Berkeley CS285 HW3 & DQN + Double DQN on LunarLander and Ms.\ Pac-Man. \href{https://rail.eecs.berkeley.edu/deeprlcourse/}{Course page} \\
\bottomrule
\end{longtable}

\subsection{Classic Environments to Implement}

\subsubsection{Cliff Walking (Example 6.6)}

4$\times$12 grid. Start at bottom-left, goal at bottom-right. The bottom row between start and goal is a ``cliff''---stepping on it gives reward $-100$ and resets to start. Every other step gives $-1$. This is THE canonical demonstration of on-policy vs off-policy behavior. SARSA accounts for its own exploration noise; Q-learning ignores it.

Implement SARSA, Q-learning, and Expected SARSA. Plot sum-of-rewards per episode for all three algorithms over 500 episodes, averaged over 100 runs.

\subsubsection{Windy Gridworld (Example 6.5)}

7$\times$10 grid with wind columns. Implement SARSA with 4 actions, then extend to King's moves (8 actions), stochastic wind. Plot timesteps vs episodes completed.

\subsubsection{Taxi-v3 (Gymnasium)}

500 discrete states. 6 actions. Compare Q-learning, SARSA, and Expected SARSA convergence rates and final win rates. \href{https://gymnasium.farama.org/environments/toy_text/taxi/}{Environment docs}.

\subsubsection{Maximization Bias / Double Q-Learning (Example 6.7)}

Implement Q-learning (observe overestimation) and Double Q-learning (observe correction). Plot \% of left actions taken from start state over episodes.

\subsection{Mini-Projects}

\subsubsection{Project 1: TD Control Benchmark}

Implement SARSA, Expected SARSA, Q-learning, and Double Q-learning from scratch (no RL libraries). Benchmark on FrozenLake-v1 (4$\times$4 and 8$\times$8, slippery=True), CliffWalking-v0, and Taxi-v3. For each algorithm $\times$ environment: run 100 independent trials, 1000 episodes each. Vary $\varepsilon \in \{0.01, 0.05, 0.1, 0.2\}$ and $\alpha \in \{0.1, 0.3, 0.5\}$. Produce heatmaps of final average return.

\subsubsection{Project 2: Safety-Aware Navigation}

Design a custom 10$\times$10 gridworld with danger zones ($-50$ but non-terminal) and cliffs ($-100$, reset). Demonstrate SARSA avoids danger zones during training while Q-learning finds the shorter but riskier path. Visualize learned Q-value heatmaps and policy arrows.

\subsection{Key Coding Resources}

\begin{itemize}[nosep]
    \item \href{https://github.com/dennybritz/reinforcement-learning}{Denny Britz SARSA/Q-learning notebooks}
    \item \href{https://github.com/SwamiKannan/CliffWalk}{Cliff Walking implementations}
    \item \href{https://github.com/georgeyiasemis/SARSA-and-Q-learning-on-a-Windy-Grid-World}{Windy Gridworld implementations}
    \item \href{https://marcinbogdanski.github.io/rl-sketchpad/RL_An_Introduction_2018/0604_Sarsa.html}{RL Sketchpad -- Ch.\ 6 figure reproductions}
\end{itemize}

\subsection{Recommended Sequence}

\begin{enumerate}[nosep]
    \item Read S\&B Ch 6.4--6.5 (SARSA, Q-learning). Do exercises 6.8, 6.11, 6.12
    \item Implement Windy Gridworld (6.6--6.7)
    \item Implement Cliff Walking---reproduce Example 6.6
    \item Implement Double Q-learning (6.14) and Taxi-v3
    \item Complete Easy21 all parts
    \item Bridge to deep RL with CS285 HW3 or CS234 Assignment 2
\end{enumerate}

\newpage
%=============================================================================
\section{Area 3: Deep Learning Foundations}
%=============================================================================

\subsection{University Assignments}

\begin{longtable}{@{}p{2in}p{3.8in}@{}}
\toprule
\textbf{Assignment} & \textbf{What You Do} \\
\midrule
\endhead
CS231n Assignment 1 & k-NN, SVM, Softmax, Two-layer net---all from scratch in NumPy on CIFAR-10. \href{https://cs231n.github.io/assignments2025/assignment1/}{Assignment page} \\
\addlinespace
CS230 Section 3 Exercises & Hand-worked backprop through computation graphs. \href{https://cs230.stanford.edu/winter2020/section3_exercises.pdf}{Exercises PDF} \\
\addlinespace
Andrew Ng's DL Specialization & Course 1 Weeks 2--4: logistic regression, 1-hidden-layer net, $L$-layer net from scratch. Course 2 Week 1: L2 reg, dropout, gradient checking. \href{https://www.coursera.org/specializations/deep-learning}{Course page} \\
\addlinespace
CS229 Problem Set 1 & Logistic regression, gradient/Hessian derivation, Newton's method vs GD. \href{https://cs229.stanford.edu/summer2020/ps1.pdf}{PDF} \\
\bottomrule
\end{longtable}

\subsection{Hand-Computation Exercises (Do on Paper)}

\begin{enumerate}[nosep]
    \item \textbf{Matt Mazur's Step-by-Step Example}: 2-input, 2-hidden, 2-output net with concrete numbers. Full forward+backward pass. \href{https://mattmazur.com/2015/03/17/a-step-by-step-backpropagation-example/}{Link}
    \item \textbf{Single-neuron logistic regression}: Derive $\partial L / \partial w$ showing the form $(\sigma(z) - y) \cdot x$
    \item \textbf{Two-layer network with MSE}: Derive all four gradient terms $\partial L / \partial W_2$, $\partial L / \partial b_2$, $\partial L / \partial W_1$, $\partial L / \partial b_1$
    \item \textbf{ReLU vs.\ Sigmoid gradient flow}: Repeat exercise 3 with ReLU. Notice why ReLU mitigates vanishing gradients
    \item \textbf{Softmax + Cross-Entropy}: Show gradient simplifies to $(\text{softmax}(z) - y_{\text{one-hot}})$
\end{enumerate}

\subsection{From-Scratch Projects}

\subsubsection{Project 1: Linear Regression from Scratch}

Pure NumPy. Implement forward pass ($\hat{y} = Xw + b$), MSE loss, gradient computation, and gradient descent update loop. Compare with the closed-form solution (normal equation). Extend to batch, stochastic, and mini-batch gradient descent.

\subsubsection{Project 2: Logistic Regression from Scratch}

Binary classifier on \texttt{make\_moons}. Implement sigmoid, binary cross-entropy, backward pass, training loop. Visualize decision boundary at different epochs. Add L2 regularization and observe changes.

\subsubsection{Project 3: Karpathy's Micrograd (Rebuild It)}

The single best exercise for understanding backpropagation and autograd. Build a scalar-valued automatic differentiation engine ($\sim$100 lines of Python). Implement a \texttt{Value} class, \texttt{backward()} using topological sort, then \texttt{Neuron}, \texttt{Layer}, and \texttt{MLP} classes.

\begin{itemize}[nosep]
    \item Repository: \href{https://github.com/karpathy/micrograd}{github.com/karpathy/micrograd}
    \item Video lecture (2+ hours): \href{https://karpathy.ai/zero-to-hero.html}{karpathy.ai/zero-to-hero.html}
\end{itemize}

\subsubsection{Project 4: Two-Layer Neural Network from Scratch}

Full NumPy implementation on MNIST. Implement parameter initialization, forward pass, cross-entropy loss, backward pass, and mini-batch training loop. Extensions: ReLU activation, L2 regularization, dropout, momentum, and Adam optimizer.

\subsubsection{Project 5: Michael Nielsen's 74-Line Network}

Work through Chapters 1--3 of \href{http://neuralnetworksanddeeplearning.com/}{neuralnetworksanddeeplearning.com}. Classifies MNIST digits at $>$96\% accuracy with no frameworks.

\subsection{Intuition-Building Mini-Projects}

\begin{itemize}[nosep]
    \item \textbf{Loss Landscape Visualization}: Visualize the loss surface along two random directions in parameter space. Reference: \href{https://arxiv.org/abs/1712.09913}{Li et al., 2017}
    \item \textbf{Gradient Flow Monitor}: Track gradient magnitudes per layer. Compare sigmoid (vanishing) vs.\ ReLU (healthy) vs.\ tanh. Add batch normalization.
    \item \textbf{Optimizer Comparison}: Implement vanilla SGD, momentum, RMSProp, Adam on the Rosenbrock function. Plot trajectories on contour plots.
    \item \textbf{Overfitting Laboratory}: Fit progressively larger networks on 20 points from a noisy sine wave. Then add L2, dropout, early stopping one at a time.
\end{itemize}

\subsection{Recommended Sequence}

\begin{enumerate}[nosep]
    \item Hand-compute backprop (Matt Mazur, CS230 exercises)
    \item Implement logistic regression from scratch
    \item Build micrograd---the most important single exercise
    \item CS231n Assignment 1 (two-layer net portion)
    \item Nielsen Chapters 1--3
\end{enumerate}

\newpage
%=============================================================================
\section{Area 4: PyTorch \& Automatic Differentiation}
%=============================================================================

\subsection{Official PyTorch Tutorials (In Order)}

\begin{enumerate}[nosep]
    \item \href{https://docs.pytorch.org/tutorials/beginner/pytorch_with_examples.html}{Learning PyTorch with Examples}
    \item \href{https://docs.pytorch.org/tutorials/beginner/blitz/autograd_tutorial.html}{Gentle Introduction to torch.autograd}
    \item \href{https://docs.pytorch.org/tutorials/beginner/introyt/autogradyt_tutorial.html}{Fundamentals of Autograd (video)}
    \item \href{https://docs.pytorch.org/tutorials/beginner/basics/optimization_tutorial.html}{Optimizing Model Parameters}
    \item \href{https://docs.pytorch.org/docs/stable/notes/autograd.html}{Autograd Mechanics Reference}
\end{enumerate}

\subsection{Build Autograd from Scratch}

\begin{longtable}{@{}p{2.2in}p{3.6in}@{}}
\toprule
\textbf{Resource} & \textbf{Description} \\
\midrule
\endhead
Karpathy's micrograd & Scalar autograd engine + MLP ($\sim$100 lines). \href{https://github.com/karpathy/micrograd}{GitHub} \\
\addlinespace
CMU 11-785 new\_grad & Tensor-valued autodiff on NumPy---closer to real PyTorch. \href{https://github.com/CMU-IDeeL/new_grad}{GitHub} \\
\addlinespace
CMU 11-785 MyTorch HW series & Build PyTorch from scratch over 4 assignments (linear layers, batch norm, conv, RNNs). \href{https://deeplearning.cs.cmu.edu/}{Course page} \\
\bottomrule
\end{longtable}

\subsection{University Assignments}

\begin{longtable}{@{}p{2.2in}p{3.6in}@{}}
\toprule
\textbf{Assignment} & \textbf{What You Do} \\
\midrule
\endhead
UMich EECS 498 Asgn.\ 4 & PyTorch Autograd and nn---three abstraction levels (raw tensors, \texttt{nn.Module}, \texttt{nn.Sequential}), culminates in ResNet. \href{https://web.eecs.umich.edu/~justincj/teaching/eecs498/FA2020/assignment4.html}{Assignment page} \\
\addlinespace
CS231n Assignment 2 & Batch normalization, dropout, conv layers from scratch, then training with PyTorch. \href{https://cs231n.github.io/assignments2025/assignment2/}{Assignment page} \\
\addlinespace
D2L.ai Ch 2.5 & Autograd exercises with PyTorch. \href{https://d2l.ai/chapter_preliminaries/autograd.html}{Link} \\
\bottomrule
\end{longtable}

\subsection{Mini-Projects}

\begin{enumerate}[nosep]
    \item \textbf{MNIST CNN}: Build with raw \texttt{nn.Conv2d}, \texttt{nn.BatchNorm2d}, \texttt{nn.Linear}. Write full training loop manually. Target $>$99\% accuracy.
    \item \textbf{ResNet-18 from Scratch on CIFAR-10}: Implement \texttt{BasicBlock} with skip connections. Compare with/without residual connections.
    \item \textbf{Normalization Comparison}: Replace BatchNorm with LayerNorm, GroupNorm, InstanceNorm in your ResNet. Compare convergence.
\end{enumerate}

\subsection{Supplementary Resources}

\begin{itemize}[nosep]
    \item \href{https://www.eletreby.me/blog/understanding-pytorch-computational-graphs}{Understanding PyTorch Computational Graphs} (blog with diagrams)
    \item \href{https://dlsys.cs.washington.edu/pdf/lecture4.pdf}{UW CSE 599W Lecture 4: Backprop and Autodiff} (systems perspective)
    \item \href{https://www.cs.toronto.edu/~rgrosse/courses/csc421_2019/readings/L06\%20Automatic\%20Differentiation.pdf}{U Toronto CSC421: Automatic Differentiation} (theoretical)
    \item \href{https://course.fast.ai/Lessons/part2.html}{fast.ai Part 2}: 30+ hours building Stable Diffusion from scratch in PyTorch
\end{itemize}

\subsection{Recommended Sequence}

\begin{enumerate}[nosep]
    \item Official PyTorch tutorials 1--5
    \item Build micrograd from scratch
    \item CMU 11-785 HW1 Part 1 (MyTorch linear layers)
    \item Mini-Projects 1 and 2 (MNIST CNN, ResNet-18)
    \item UMich EECS 498 Assignment 4
\end{enumerate}

\newpage
%=============================================================================
\section{Area 5: DAVI and DQNs}
%=============================================================================

\subsection{Papers to Read (In Order)}

\begin{enumerate}[nosep]
    \item \textbf{Mnih et al.\ (2013)}: Original DQN. \href{https://arxiv.org/abs/1312.5602}{arXiv:1312.5602}
    \item \textbf{Mnih et al.\ (2015)}: Nature DQN (adds target networks). \href{https://www.nature.com/articles/nature14236}{Nature}
    \item \textbf{Agostinelli et al.\ (2019)}: DeepCubeA---uses DAVI + batch weighted A* search. \href{https://www.nature.com/articles/s42256-019-0070-z}{Nature MI}. Code: \href{https://github.com/forestagostinelli/DeepCubeA}{GitHub}
    \item \textbf{van Hasselt et al.\ (2016)}: Double DQN. \href{https://arxiv.org/abs/1509.06461}{arXiv:1509.06461}
    \item \textbf{Schaul et al.\ (2016)}: Prioritized Experience Replay. \href{https://arxiv.org/abs/1511.05952}{arXiv:1511.05952}
    \item \textbf{van Hasselt et al.\ (2018)}: Deadly Triad. \href{https://arxiv.org/abs/1812.02648}{arXiv:1812.02648}
\end{enumerate}

\subsection{University Assignments}

\begin{longtable}{@{}p{2.2in}p{3.6in}@{}}
\toprule
\textbf{Assignment} & \textbf{What You Do} \\
\midrule
\endhead
Stanford CS234 Asgn.\ 2 & Linear Q-learning to full DQN on Atari (experience replay + target networks). \href{https://web.stanford.edu/class/cs234/}{Course page} \\
\addlinespace
Berkeley CS285 HW3 & DQN + Double DQN on LunarLander and Ms.\ Pac-Man. \href{https://rail.eecs.berkeley.edu/deeprlcourse/}{Course page} \\
\addlinespace
CMU 10-703 HW2 & Deep Q-networks with online + target networks. \href{https://cmudeeprl.github.io/703website_f24/}{Course page} \\
\bottomrule
\end{longtable}

\subsection{Progressive Implementation Path}

\begin{longtable}{@{}clp{3.2in}@{}}
\toprule
\textbf{Level} & \textbf{Project} & \textbf{Environment} \\
\midrule
\endhead
1 & Tabular Q-learning & FrozenLake-v1 \\
2 & Linear Q-learning (observe instability) & CartPole-v1 \\
3 & Full DQN (replay, target net, $\varepsilon$-greedy) & CartPole-v1 \\
4 & DQN with tuning & LunarLander-v2 \\
5 & DQN on pixels (CNN, frame stacking) & Pong / Breakout \\
6 & Double DQN + Prioritized Replay & Atari \\
7 & \textbf{DAVI on combinatorial puzzle} & 15-puzzle or 2$\times$2 Rubik's cube \\
\bottomrule
\end{longtable}

\subsection{Key Concepts}

\subsubsection{Target Networks}

If the same network generates both the prediction and the target, every update shifts the target, creating a feedback loop. The target network $\theta^{-}$ freezes the target for $N$ steps. The DQN loss becomes:
\begin{align}
    L(\theta) = \mathbb{E}\left[\left(r + \gamma \max_{a'} Q(s', a'; \theta^{-}) - Q(s, a; \theta)\right)^2\right]
\end{align}

\subsubsection{Experience Replay}

Online RL training data is (a) temporally correlated and (b) non-stationary. Both properties violate the i.i.d.\ assumption of SGD. Experience replay stores transitions $(s, a, r, s')$ in a buffer and samples mini-batches uniformly at random, breaking both correlations.

\subsubsection{The Deadly Triad}

Function approximation + bootstrapping + off-policy learning can cause divergence. DQN mitigates this with target networks and experience replay, but does not fully solve it.

\subsubsection{Deep Approximate Value Iteration (DAVI)}

Generate training data by scrambling from the goal state. Train a neural network to predict cost-to-go. Use that network as a heuristic in A* search. This is the approach in Prof.\ Agostinelli's DeepCubeA paper.

\subsection{Key Implementation References}

\begin{itemize}[nosep]
    \item \href{https://github.com/vwxyzjn/cleanrl/blob/master/cleanrl/dqn.py}{CleanRL DQN}---single-file, most readable DQN
    \item \href{https://docs.pytorch.org/tutorials/intermediate/reinforcement_q_learning.html}{PyTorch DQN Tutorial}---CartPole from scratch
    \item \href{https://huggingface.co/learn/deep-rl-course/en/unit3/introduction}{Hugging Face Deep RL Course Unit 3}---theory + hands-on notebook
    \item \href{https://araffin.github.io/post/rl102/}{Antonin Raffin: Tabular to DQN}---best progressive walkthrough
    \item \href{https://danieltakeshi.github.io/2016/11/25/frame-skipping-and-preprocessing-for-deep-q-networks-on-atari-2600-games/}{Daniel Takeshi: Atari Preprocessing}
\end{itemize}

\subsection{Recommended Sequence}

\begin{enumerate}[nosep]
    \item Read Sutton \& Barto Ch.\ 9--11 (function approximation theory)
    \item Read DQN papers (2013, then 2015 Nature)
    \item Implement tabular Q-learning on FrozenLake, then linear FA on CartPole
    \item Implement full DQN on CartPole, then LunarLander
    \item Read DeepCubeA paper, study the \href{https://github.com/forestagostinelli/DeepCubeA}{GitHub repo}
    \item Implement simplified DAVI on a combinatorial puzzle
\end{enumerate}

\newpage
%=============================================================================
\section{Solution References for Self-Checking}
%=============================================================================

\begin{longtable}{@{}p{2.6in}p{3.2in}@{}}
\toprule
\textbf{Resource} & \textbf{Link} \\
\midrule
\endhead
Sutton \& Barto Solutions (LyWangPX) & \href{https://github.com/LyWangPX/Reinforcement-Learning-2nd-Edition-by-Sutton-Exercise-Solutions}{GitHub} \\
\addlinespace
Tianlin Liu's Solutions PDF & \href{https://tianlinliu.com/files/notes_exercise_RL.pdf}{PDF} \\
\addlinespace
Scott Jeen's Exercise Compilation & \href{https://enjeeneer.io/sutton_and_barto/rl_exercises.pdf}{PDF} \\
\addlinespace
Easy21 Solutions (alexandrasouly) & \href{https://github.com/alexandrasouly/easy21}{GitHub} \\
\addlinespace
CS234 Solutions (ksang) & \href{https://github.com/ksang/cs234-assignments}{GitHub} \\
\addlinespace
CS285 Solutions (ZHZisZZ) & \href{https://github.com/ZHZisZZ/cs285-homework-fall2021}{GitHub} \\
\addlinespace
Denny Britz RL Repository & \href{https://github.com/dennybritz/reinforcement-learning}{GitHub} \\
\addlinespace
Shangtong Zhang Reproductions & \href{https://github.com/ShangtongZhang/reinforcement-learning-an-introduction}{GitHub} \\
\bottomrule
\end{longtable}

\vspace{12pt}
\hrule height 1pt
\vspace{6pt}
\begin{center}
\textit{The single assignment that ties Areas 1 and 2 together most effectively is \textbf{Easy21}. For Areas 3 and 4, \textbf{micrograd} is the highest-value single exercise. For Area 5, start with CleanRL's DQN and work up to studying Prof.\ Agostinelli's DeepCubeA code.}
\end{center}

\end{document}
